\chapter{Formalization map}
\label{ch:formal-map}

\section{The dependency graph}

The intended Lean dependency graph is:
\[
\begin{array}{c}
\text{complex analysis and topology}\\
\downarrow\\
\text{quadratic dynamics, Green functions, Böttcher coordinates}\\
\downarrow\\
\text{puzzles and parapuzzles}\\
\downarrow\\
\text{modulus, a priori bounds, renormalization}\\
\downarrow\\
\text{phase--parameter correspondence}\\
\downarrow\\
\text{local connectedness at a parameter class}\\
\downarrow\\
\text{global MLC}.
\end{array}
\]
Each downward arrow is a separate proof obligation; the diagram is
an intended mathematical architecture.  A result may be
imported from an existing Lean dependency only when its exact hypotheses and
conclusion match the statement used here.

\section{Current Lean status}

\begin{center}\small
\begin{tabularx}{0.98\textwidth}{@{}P{0.25\textwidth}@{\hspace{8pt}}Y@{\hspace{8pt}}P{0.13\textwidth}@{}}
\textbf{Layer} & \textbf{Current content} & \textbf{Status}\\
\hline
Orbit and outer approximations &
Escape criterion, compact finite stages, and exact intersection. &
proved\\
Periods one and two &
Multiplier product, cardioid and period-two parametrizations, frontiers,
uniqueness of attracting fixed points, disjointness, and the common root $-3/4$. &
proved\\
Real slice of $\M$ &
$\M\cap\R=[-2,1/4]$, with escape and invariant-interval estimates. &
proved\\
Multiplier map on the period-two disk &
affine biholomorphism $c\mapsto4(c+1)$ onto $\mathbb D$. &
proved\\
Multiplier map on $\mathcal H_0$ &
biholomorphism onto $\mathbb D$ with explicit branch $1-\sqrt{1-4c}$. &
proved\\
Real parameters of $\mathcal H_0$ and $\mathcal H_2$ &
$\mathcal H_0\cap\R=(-\tfrac34,\tfrac14)$ and $\mathcal H_2\cap\R=(-\tfrac54,-\tfrac34)$. &
proved\\
Exact examples for $c=0$ and $c=-2$ &
$K(0)$ is the closed disk and $J(0)$ the circle; $K(-2)=J(-2)=[-2,2]$. &
proved\\
Local component criterion &
Uniform outer buffer implies local connectedness. &
proved conditionally\\
Frozen modulus audit &
Gaussian nested differences are summable; the legacy finite class is empty. &
proved obstruction\\
Classical Yoccoz shrinking &
Graph cuts and conformal-modulus estimates. &
pending\\
Parameter parapuzzles &
Connected induced pieces and a basis interface. &
partial\\
Phase--parameter contracts &
External rays, dynamical graphs, parameter graphs, and topological transport. &
interface\\
Categorical branch &
TopCat over-approximations, conditional roots, carving interfaces, and the
refuted frozen Green-sublevel proposal. &
mixed\\
Renormalization towers &
Discrete indices and frozen weighted-area contracts; positive lower-bound
certificate refuted. &
model audit\\
Bounded-type satellite and Feigenbaum cases &
Named literature interfaces and source ledger. &
open formalization\\
Parabolically bounded primitive renormalization &
Thin--Thick, Value Calculus, Wanderers, and Wave Lemma inputs. &
open formalization\\
Neutral renormalization &
Bubble bounds, attractor rigidity, and hybrid-class geometry. &
open formalization\\
Dudko--Lyubich/KKL case certificates &
Source-specific class predicates plus connected shrinking parameter pieces. &
interface\\
Combinatorial profiles &
Coarse type labels, additive depth, multiplicative return period,
and nonnegative-series arithmetic. &
proved arithmetic\\
Global MLC &
Target \lean{LocallyConnectedSpace M}. &
open\\
\end{tabularx}
\end{center}

The module \lean{Mlc.Literature.Review} is deliberately conservative:
\begin{itemize}
\item \lean{yoccoz_puzzle_shrinkage} is a theorem of the frozen
      model with an impossible premise;
\item \lean{FinitelyRenormalizableMLC} concerns the same empty
      legacy class;
\item \lean{BoundedTypeSatelliteMLC} and \lean{FeigenbaumMLC}
      are parameterized targets whose actual source-class
      definitions still have to be implemented.
\end{itemize}
No external paper is silently converted into an axiom.

The legacy structure \lean{ParaPuzzleCertificate} is tied to the
frozen full-sublevel pieces.  Its connectedness and singleton-intersection fields
imply local connectivity through the theorem
\begin{center}
\lean{localConnectivity_of_paraPuzzleCertificate}.
\end{center}
Its corresponding family proposition is
\begin{center}
\lean{YoccozParameterPuzzleResult},
\end{center}
quantified over every $c\in\M$.  Unlike the finite family in
Chapter~\ref{chap:theorem-matrix}, it is not indexed by the empty
legacy finite class.
Its frozen-piece existence requirement is separate from the
classical graph-cut construction and cannot be supplied merely by
citing a Yoccoz theorem.

The new module \lean{Mlc.Literature.Foundations} gives typed names for the
quadratic map, orbit, filled Julia set, Julia set, and Mandelbrot set.  Its
main structural interfaces are:
\begin{itemize}
\item \lean{ParameterPuzzleBasis};
\item \lean{APrioriBounds};
\item \lean{RenormalizationMLCInterface}.
\end{itemize}
These records specify the shape of required data, rather than
constructing those data or identifying an arbitrary real-valued
field with conformal modulus.
The constructor
\leanref{Mlc/Literature/Foundations.lean}{ParameterPuzzleBasis.of_compactClosure}
supplies the topological basis from actual compact parameter pieces.
The audit in
\leanref{Mlc/Literature/ModelAudit.lean}{summable_weightedArea_differences}
separately explains why the inherited analytic model must be replaced.

\section{The current conditional root}

The existing root is
\[
\mathrm{BUF}(\M,O)\Longrightarrow\mathrm{MLC}(\M),
\]
implemented by
\lean{MLC.mlc_conjecture_of_uniformOuterBuffer}.  The alternative uniform
geometric route and the finite-component route are useful independent
interfaces, but their existence hypotheses remain open.  The review keeps
these routes visible because they are possible formalization targets, not
because they have already proved MLC.
